\documentclass[11pt,a4paper]{article}
\input{../template/utm-lab}

\newcommand{\labnumber}{6}
\newcommand{\labtitle}{Inside the xv6 File System}
\newcommand{\studentname}{Daniel Nicorici}
\newcommand{\studentgroup}{FAF-243}

\begin{document}

\utmtitlepage
\tableofcontents
\newpage

\plainsection{Purpose of the work}

The xv6 file system uses the same ideas as ext4 (inodes, directories stored as tables of
entries, hard links and a journal), but in a few hundred readable lines. The purpose of this lab
is to read what xv6 shows about its files, to write a \code{mystat} program that prints an
inode's fields and use it to watch hard links at work, to measure the limits defined in
\code{kernel/fs.h}, and to write \code{find}, which walks the directory tree recursively.

The work was done in my \code{\textasciitilde/xv6-riscv} tree (\code{mit-pdos/xv6-riscv}, commit
\code{06aad25}), cross-compiled with \code{riscv64-linux-gnu-gcc} 16.2.1 and run on QEMU 11.1.1
on my Arch Linux laptop. Every experiment was run on a fresh disk image. Each new program is a
file \code{user/NAME.c} plus a line \code{\$U/\_NAME\textbackslash} in \code{UPROGS}.

\section{Part 1. Read what xv6 shows you}

\begin{lstlisting}[caption={ls on a fresh xv6 disk}]
$ ls
.              1 1 1024
..             1 1 1024
README         2 2 2441
cat            2 3 35904
echo           2 4 34800
forktest       2 5 17704
grep           2 6 43328
init           2 7 35144
kill           2 8 34696
ln             2 9 34496
ls             2 10 41984
mkdir          2 11 34752
rm             2 12 34728
sh             2 13 57104
stressfs       2 14 35576
usertests      2 15 201552
grind          2 16 50376
wc             2 17 36872
zombie         2 18 34048
logstress      2 19 36720
forphan        2 20 35512
dorphan        2 21 34968
sync           2 22 34144
sleep          2 23 34472
args           2 24 34464
mystat         2 25 35424
find           2 26 41336
console        3 27 0
\end{lstlisting}

Each line is: name, type, inode number and size in bytes. My disk also contains the programs
from earlier labs (\code{sleep}, \code{args}) and this lab (\code{mystat}, \code{find}), so
\code{console} has inode 27 instead of the 25 shown in the assignment.

\begin{lstlisting}[caption={An empty directory}]
$ mkdir a
$ ls a
.              1 28 32
..             1 1 1024
\end{lstlisting}

\begin{lstlisting}[style=ccode,caption={File types and struct stat in kernel/stat.h}]
#define T_DIR    1 // Directory
#define T_FILE   2 // File
#define T_DEVICE 3 // Device

struct stat {
  int dev;     // File system's disk device
  uint ino;    // Inode number
  short type;  // Type of file
  short nlink; // Number of links to file
  uint64 size; // Size of file in bytes
};
\end{lstlisting}

\begin{lstlisting}[style=ccode,firstnumber=54,caption={A directory entry in kernel/fs.h}]
#define DIRSIZ 14

// The name field may have DIRSIZ characters and not end in a NUL
// character.
struct dirent {
  ushort inum;
  char name[DIRSIZ] __attribute__((nonstring));
};
\end{lstlisting}

\textbf{Observations:}
\begin{enumerate}
  \item \textbf{Types 1, 2 and 3.} From \code{kernel/stat.h}: 1 is \code{T\_DIR}, a directory;
  2 is \code{T\_FILE}, a regular file; 3 is \code{T\_DEVICE}, a device. \code{console} is the
  only device, the terminal, exactly like \code{/dev/tty} on Linux.
  \item \textbf{The root directory.} The root has inode number 1. Both \code{.} and
  \code{..} show inode 1 because the root has no parent: when \code{mkfs} creates the disk, it
  makes the root's \code{..} point back to the root itself, so going ``up'' from \code{/}
  stays in \code{/}.
  \item \textbf{Size of a directory entry.} A directory is just a file filled with
  \code{struct dirent} records: a 2-byte inode number (\code{ushort inum}) and a 14-byte name
  (\code{DIRSIZ}), 16 bytes in total. The empty directory \code{a} has size 32, which is two
  entries: \code{.} pointing to \code{a} itself (inode 28) and \code{..} pointing to its
  parent, the root (inode 1).
\end{enumerate}

\section{Part 2. Write mystat}

\code{stat()} asks the kernel to fill a \code{struct stat} for a path. The program converts
the type number into a word with a \code{switch}, prints an error to file descriptor 2 for a
name that does not exist and continues with the next one, and prints the four fields for every
name that does.

\begin{lstlisting}[style=ccode,caption={user/mystat.c}]
#include "kernel/types.h"
#include "kernel/stat.h"
#include "user/user.h"

char *
typename(int t)
{
  switch(t){
  case T_DIR:
    return "directory";
  case T_FILE:
    return "file";
  case T_DEVICE:
    return "device";
  }
  return "?";
}

int
main(int argc, char *argv[])
{
  struct stat st;
  for(int i = 1; i < argc; i++){
    if(stat(argv[i], &st) < 0){
      fprintf(2, "mystat: cannot stat %s\n", argv[i]);
      continue;
    }
    printf("%s: %s, inode %d, links %d, size %d bytes\n",
           argv[i], typename(st.type), st.ino, st.nlink, (int)st.size);
  }
  exit(0);
}
\end{lstlisting}

\begin{lstlisting}[caption={The UPROGS lines for mystat and find}]
$ grep -n "_mystat\|_find" Makefile
155:	$U/_mystat\
156:	$U/_find\
\end{lstlisting}

\begin{lstlisting}[caption={mystat watching a hard link}]
$ mystat README . console
README: file, inode 2, links 1, size 2441 bytes
.: directory, inode 1, links 1, size 1024 bytes
console: device, inode 27, links 1, size 0 bytes
$ ln README r2
$ mystat README r2
README: file, inode 2, links 2, size 2441 bytes
r2: file, inode 2, links 2, size 2441 bytes
$ rm README
$ mystat r2
r2: file, inode 2, links 1, size 2441 bytes
$ cat r2 | wc
48 336 2441
$ mystat README nosuchfile
mystat: cannot stat README
mystat: cannot stat nosuchfile
\end{lstlisting}

\begin{lstlisting}[style=ccode,firstnumber=235,caption={sys\_unlink in kernel/sysfile.c}]
  memset(&de, 0, sizeof(de));
  if (writei(dp, 0, (uint64)&de, off, sizeof(de)) != sizeof(de))
    panic("unlink: writei");
  if (ip->type == T_DIR) {
    dp->nlink--;
    iupdate(dp);
  }
  iunlockput(dp);

  ip->nlink--;
  iupdate(ip);
  iunlockput(ip);
\end{lstlisting}

\begin{lstlisting}[style=ccode,firstnumber=357,caption={iput in kernel/fs.c, where the blocks are freed}]
  int last = (ip->ref == 1 && ip->valid && ip->nlink == 0);
  uint dev = ip->dev, inum = ip->inum;

  if (last) {
    // ip->ref == 1 means no other process can have ip locked.
    acquiresleep(&ip->lock);
    release(&itable.lock);

    itrunc(ip); // free the data blocks (type stays nonzero on disk)
    ip->valid = 0;

    releasesleep(&ip->lock);

    acquire(&itable.lock);
  }

  ip->ref--;
  release(&itable.lock);

  if (last)
    ifree(dev, inum); // now clear type on disk: inum becomes allocatable
\end{lstlisting}

\begin{lstlisting}[style=ccode,firstnumber=32,caption={The on-disk inode, struct dinode in kernel/fs.h}]
struct dinode {
  short type;              // File type
  short major;             // Major device number (T_DEVICE only)
  short minor;             // Minor device number (T_DEVICE only)
  short nlink;             // Number of links to inode in file system
  uint size;               // Size of file (bytes)
  uint addrs[NDIRECT + 1]; // Data block addresses
};
\end{lstlisting}

\textbf{Observations:}
\begin{enumerate}
  \item \textbf{After \code{ln} and after \code{rm}.} \code{ln README r2} did not copy
  anything: both names show the same inode 2, the same size, and the link count went from 1
  to 2. After \code{rm README}, the name \code{README} is gone (\code{mystat} can no longer
  find it), but \code{r2} still shows inode 2 with the full 2441 bytes, and \code{wc} confirms
  the content is intact; only the link count dropped back to 1. What changed is the set of
  names and the link count; what stayed the same is the inode number, the size and the data.
  \item \textbf{Where the link count is decremented, and when the blocks are freed.} In
  \code{sys\_unlink} (\code{kernel/sysfile.c}), line 236 overwrites the directory entry with
  zeros, which removes the name, and line 244, \code{ip->nlink--}, decrements the inode's link
  count. The data blocks are not freed there. They are freed in \code{iput()}
  (\code{kernel/fs.c}): only when the link count is 0 \emph{and} this is the last in-memory
  reference to the inode (\code{ref == 1}, so no process still has the file open) does
  \code{itrunc()} on line 365 free the data blocks, after which \code{ifree()} marks the inode
  itself as free. So deleting a file that some program still has open keeps its data until the
  file is closed.
  \item \textbf{Fields of an xv6 inode.} \code{struct dinode} has only \code{type},
  \code{major} and \code{minor} (device numbers, used only for devices), \code{nlink},
  \code{size} and \code{addrs[NDIRECT + 1]}, the addresses of the data blocks. There is no
  owner, no group, no permission bits and no timestamps, and no inode type for symbolic links,
  which is why xv6 has neither permissions nor \code{ln -s}.
\end{enumerate}

\section{Part 3. Measure the limits}

\begin{lstlisting}[caption={The limits defined in kernel/fs.h}]
$ grep -n "define BSIZE\|define NDIRECT\|define NINDIRECT\|define DIRSIZ" kernel/fs.h
5:#define BSIZE   1024 // block size
26:#define NDIRECT   12
27:#define NINDIRECT (BSIZE / sizeof(uint))
54:#define DIRSIZ 14
\end{lstlisting}

\begin{lstlisting}[caption={The disk layout printed by mkfs}]
$ make clean ; make fs.img 2>&1 | grep nmeta
nmeta 47 (boot, super, log blocks 31, inode blocks 13, bitmap blocks 1) blocks 1953 total 2000
\end{lstlisting}

\begin{lstlisting}[caption={A name longer than DIRSIZ}]
$ mkdir averyveryverylongname
$ ls
.              1 1 1024
..             1 1 1024
README         2 2 2441
[... 25 entries omitted ...]
a              1 28 32
averyveryveryl 1 29 32
\end{lstlisting}

\textbf{Observations:}
\begin{enumerate}
  \item \textbf{The long name.} \code{ls} shows \code{averyveryveryl}: only the first 14
  characters were kept. The constant is \code{DIRSIZ = 14}, the size of the name field in
  \code{struct dirent}; \code{dirlink()} copies the name with
  \code{strncpy(de.name, name, DIRSIZ)} (\code{kernel/fs.c}, line 640), which silently cuts
  everything after 14 characters.
  \item \textbf{The largest file.} An inode has \code{NDIRECT = 12} direct block addresses plus
  one indirect block. The indirect block is a block of \code{BSIZE = 1024} bytes full of 4-byte
  block numbers, so it holds \code{NINDIRECT} \(= 1024 / 4 = 256\) more addresses. A file can
  therefore use at most \(12 + 256 = 268\) blocks (\code{MAXFILE}), and the largest file is
  \[ 268 \times 1024 = 274\,432 \text{ bytes} \approx 268 \text{ KB}. \]
  \item \textbf{The disk layout.} The whole disk has 2000 blocks of 1 KB
  (\code{FSSIZE = 2000}). 47 of them are metadata: 1 boot block, 1 superblock, 31 log blocks
  (a header plus \code{LOGBLOCKS = 30} for the journal), 13 inode blocks and 1 bitmap block. The
  remaining 1953 blocks hold data. The 13 inode blocks follow from the sizes: a
  \code{struct dinode} is 64 bytes, so 16 fit in a 1024-byte block, and
  \(200 / 16 + 1 = 13\) blocks hold the 200 inodes that \code{mkfs} creates.
\end{enumerate}

\section{Part 4. Write find}

\code{find} starts from \code{user/ls.c}: it opens a directory, checks with \code{fstat} that
it really is one, and reads it as a sequence of 16-byte \code{struct dirent} records. For every
used entry it builds the full path in a buffer (\code{path/name}), calls \code{stat} on it,
prints the path if the last component equals the name being searched, and calls itself
recursively if the entry is a directory. Entries with \code{inum == 0} are free slots (for
example left by \code{rm}) and are skipped, and so are \code{.} and \code{..}, otherwise the
walk would loop forever. Every directory is closed before \code{find} returns, because a
process may have only 16 open files (\code{NOFILE}) and the recursion keeps one directory open
per level. The optional \code{-n} argument adds the inode number to each line.

\begin{lstlisting}[style=ccode,caption={user/find.c}]
#include "kernel/types.h"
#include "kernel/stat.h"
#include "user/user.h"
#include "kernel/fs.h"
#include "kernel/fcntl.h"

int shownum;

void
find(char *path, char *name)
{
  char buf[512], *p;
  int fd;
  struct dirent de;
  struct stat st;

  if((fd = open(path, O_RDONLY)) < 0){
    fprintf(2, "find: cannot open %s\n", path);
    return;
  }
  if(fstat(fd, &st) < 0 || st.type != T_DIR){
    fprintf(2, "find: %s is not a directory\n", path);
    close(fd);
    return;
  }
  if(strlen(path) + 1 + DIRSIZ + 1 > sizeof buf){
    fprintf(2, "find: path too long\n");
    close(fd);
    return;
  }

  strcpy(buf, path);
  p = buf + strlen(buf);
  *p++ = '/';
  while(read(fd, &de, sizeof(de)) == sizeof(de)){
    if(de.inum == 0)
      continue;
    memmove(p, de.name, DIRSIZ);
    p[DIRSIZ] = 0;
    if(strcmp(p, ".") == 0 || strcmp(p, "..") == 0)
      continue;
    if(stat(buf, &st) < 0){
      fprintf(2, "find: cannot stat %s\n", buf);
      continue;
    }
    if(strcmp(p, name) == 0){
      if(shownum)
        printf("%s %d\n", buf, st.ino);
      else
        printf("%s\n", buf);
    }
    if(st.type == T_DIR)
      find(buf, name);
  }
  close(fd);
}

int
main(int argc, char *argv[])
{
  if(argc == 4 && strcmp(argv[3], "-n") == 0)
    shownum = 1;
  else if(argc != 3){
    fprintf(2, "usage: find <dir> <name> [-n]\n");
    exit(1);
  }
  find(argv[1], argv[2]);
  exit(0);
}
\end{lstlisting}

\begin{lstlisting}[caption={Testing find on a fresh disk}]
$ mkdir a
$ mkdir a/b
$ echo x > a/b/r2
$ echo y > a/r2
$ ln cat r2
$ find . r2
./a/b/r2
./a/r2
./r2
$ find a b
a/b
\end{lstlisting}

\begin{lstlisting}[caption={Stretch: inode numbers with -n, and the error cases}]
$ find . r2 -n
./a/b/r2 30
./a/r2 31
./r2 3
$ mystat cat r2 a/r2 a/b/r2
cat: file, inode 3, links 2, size 35904 bytes
r2: file, inode 3, links 2, size 35904 bytes
a/r2: file, inode 31, links 1, size 2 bytes
a/b/r2: file, inode 30, links 1, size 2 bytes
$ find
usage: find <dir> <name> [-n]
$ find README r2
find: README is not a directory
\end{lstlisting}

\textbf{Results:}
\begin{enumerate}
  \item \textbf{The test.} \code{find . r2} printed the three expected paths and
  \code{find a b} printed \code{a/b}, exactly as in the assignment.
  \item \textbf{Stretch.} With \code{-n}, \code{./r2} has inode 3, the same inode as
  \code{cat}, and \code{mystat} shows both names with 2 links and the size of the \code{cat}
  program: \code{ln cat r2} made \code{./r2} a second name for the \code{cat} file. The other
  two \code{r2} files (inodes 30 and 31) are separate 2-byte files created with \code{echo}.
\end{enumerate}

\plainsection{Conclusions}

This lab showed the xv6 file system from the inside. A file is an inode, a fixed-size record
that holds its type, link count, size and the addresses of its data blocks, while a name lives
only in a directory, which is itself an ordinary file filled with 16-byte entries that pair a
name with an inode number. Because of this split, a hard link is just one more entry pointing to
the same inode, and removing a name only decrements the link count. The \code{mystat}
experiment made this visible, and the kernel code confirmed that the data blocks are freed only
when no name and no open file refers to the inode any more.

Reading \code{kernel/fs.h} showed how small and fixed the design is: 1 KB blocks, 14-character
names, 12 direct and 256 indirect block addresses, which limit a file to 274\,432 bytes, and a
2000-block disk of which 47 blocks are metadata. Writing \code{find} combined all of this:
reading a directory is reading its entries one by one, every entry is turned into a path and
inspected with \code{stat}, and the tree is walked by recursion, while taking care to skip free
slots, \code{.} and \code{..}, and to close every directory because a process can hold only 16
open files. The same ideas, with more fields and bigger limits, are what ext4 does on Linux.

\end{document}
